\chapter{Requirements Specification}

\section{Functional requirements}

\subsection{Asset selection}
\begin{tabularx}{\textwidth}{@{}lX@{}}
\toprule
\textbf{ID} & \textbf{Requirement} \\ \midrule
FR-SEL-01 & Read COTAHIST files downloaded manually from the B3 website, without automated collection. \\
FR-SEL-02 & Keep only stocks and units from the cash market in standard lots. \\
FR-SEL-03 & At each rebalancing date, rank assets by average daily traded value over the 252 previous sessions, excluding assets with presence below 80\%. \\
FR-SEL-04 & Keep one share class per company and require enough history to estimate $\Sigma$. \\
FR-SEL-05 & Select the 10 most liquid assets for the study portfolio. \\
FR-SEL-06 & Generate the application catalog with the 30 most liquid assets from the most recent ranking, without publishing prices or volumes. \\
\bottomrule
\end{tabularx}

\subsection{Data collection}
\begin{tabularx}{\textwidth}{@{}lX@{}}
\toprule
\textbf{ID} & \textbf{Requirement} \\ \midrule
FR-COL-01 & Download adjusted closing prices from a free source, initially \texttt{yfinance}, for each asset in the universe. \\
FR-COL-02 & Append the \texttt{.SA} suffix to tickers and validate that they exist; an invalid ticker produces an error message. \\
FR-COL-03 & Collect the daily CDI rate from the Central Bank of Brazil API (SGS, series 12). \\
FR-COL-04 & Collect the current market capitalization of each asset for the application prior. \\
FR-COL-05 & Collect historical share counts from CVM: FRE since 2010 and the ITR capital composition since 2020, using the delivery date of each document. \\
FR-COL-06 & Accept external datasets in a fixed file format, so that the backtest does not depend on the data source. \\
FR-COL-07 & Run a routine outside the server on business days to generate a snapshot with $\Sigma$, market capitalizations, $\delta$ and the previous close of the catalog assets. \\
FR-COL-08 & If the routine fails, keep the last valid snapshot and show the date of the data to the user. \\
\bottomrule
\end{tabularx}

\subsection{Data treatment and estimation}
\begin{tabularx}{\textwidth}{@{}lX@{}}
\toprule
\textbf{ID} & \textbf{Requirement} \\ \midrule
FR-TRA-01 & Align all series on the ANBIMA business-day calendar; days with CDI and no price follow the missing data rule. \\
FR-TRA-02 & Fill missing values with forward-fill of up to 5 business days; exclude, with a warning, assets with insufficient coverage. \\
FR-TRA-03 & Compute simple returns $R_t = P_t/P_{t-1} - 1$. \\
FR-TRA-04 & Compute excess returns over the CDI: $r^{ex} = r - r_f$. \\
FR-TRA-05 & Estimate the annualized sample covariance matrix $\Sigma$ of excess returns ($\times 252$ for daily data and $\times 52$ for weekly data) and check that it is positive semidefinite. \\
\bottomrule
\end{tabularx}

\subsection{Black-Litterman model}
\begin{tabularx}{\textwidth}{@{}lX@{}}
\toprule
\textbf{ID} & \textbf{Requirement} \\ \midrule
FR-MOD-01 & Compute market weights $w_{mkt,i} = MC_i / \sum_j MC_j$ within the universe. \\
FR-MOD-02 & Compute the prior by reverse optimization: $\Pi = \delta \Sigma w_{mkt}$. \\
FR-MOD-03 & Allow $\delta$ to be set by the user or estimated as $\delta = (E[r_m]-r_f)/\sigma_m^2$. \\
FR-MOD-04 & Allow absolute views: ``asset $i$ will return $q$\% p.a. above the CDI''. \\
FR-MOD-05 & Allow relative views: ``asset $i$ will outperform asset $j$ by $q$\% p.a.''. \\
FR-MOD-06 & Automatically build and validate the $P$ and $Q$ matrices from the views. \\
FR-MOD-07 & Build $\Omega = \mathrm{diag}(P\tau\Sigma P^{\top})$, with $\tau = 0.05$ by default. \\
FR-MOD-08 & Compute the posterior $\mu_{BL} = \Pi + \tau\Sigma P^{\top}(P\tau\Sigma P^{\top} + \Omega)^{-1}(Q - P\Pi)$, solving linear systems without explicit matrix inversion. \\
FR-MOD-09 & With no views, the posterior must equal the prior. \\
FR-MOD-10 & Compute optimal weights by maximizing $w^{\top}\mu_{BL} - \frac{\delta}{2} w^{\top}\Sigma w$, with weights that sum to 1, no short selling and a maximum weight of 20\% per asset. \\
FR-MOD-11 & Display the weights table prior $\times$ BL $\times$ difference. \\
FR-MOD-12 & Compute the expected return, volatility and Sharpe ratio of the portfolio. \\
FR-MOD-13 & Report infeasible optimization problems with a clear message. \\
\bottomrule
\end{tabularx}

\subsection{Backtest study}
\begin{tabularx}{\textwidth}{@{}lX@{}}
\toprule
\textbf{ID} & \textbf{Requirement} \\ \midrule
FR-BKT-01 & Run a walk-forward backtest without views: at each rebalancing date, use only data available before that date (no look-ahead). \\
FR-BKT-02 & Use weekly data, rebalance on the first session of January, May and September, and estimate $\Sigma$ with a 3-year window, using 5 years as a robustness check. \\
FR-BKT-03 & Compare against the benchmarks: Ibovespa, CDI, the equal-weight portfolio (1/N) and the market capitalization portfolio. \\
FR-BKT-04 & Compute cumulative return, annualized volatility, Sharpe ratio and maximum drawdown. \\
FR-BKT-05 & Plot the equity curve (base 100). \\
FR-BKT-06 & Compute historical market capitalization from unadjusted prices and the latest share count delivered before the rebalancing date, keeping both on the same split basis. \\
FR-BKT-07 & If an asset stops trading, sell the position at the last price and keep the value in CDI until the next rebalancing; report results with and without reconstructed delisted assets. \\
FR-BKT-08 & Apply transaction costs to turnover, with a configurable cost per side. \\
FR-BKT-09 & (Optional) Run exercises with views separately from the main backtest and identify them clearly. \\
\bottomrule
\end{tabularx}

\subsection{API}
\begin{tabularx}{\textwidth}{@{}lX@{}}
\toprule
\textbf{ID} & \textbf{Requirement} \\ \midrule
FR-API-01 & Provide the catalog of assets. \\
FR-API-02 & Receive assets, parameters and views, and return weights, metrics and the parameters used. \\
FR-API-03 & Provide the documented results of the backtest study. \\
FR-API-04 & Return errors in a standard format with clear messages. \\
\bottomrule
\end{tabularx}

\subsection{User interface}
\begin{tabularx}{\textwidth}{@{}lX@{}}
\toprule
\textbf{ID} & \textbf{Requirement} \\ \midrule
FR-UI-01 & Start with the 10 most liquid assets and allow the user to edit the selection within the catalog, keeping at least 5 assets. \\
FR-UI-02 & Allow the user to select an asset and create an absolute or relative view. \\
FR-UI-03 & Display the weights, metrics and study backtest results. \\
FR-UI-04 & (Optional) Track the finished portfolio with the previous business day's close stored in the snapshot, showing the date of the data. \\
\bottomrule
\end{tabularx}

\section{Non-functional requirements}
\begin{itemize}
  \item Backend implemented in Python and deployed on PythonAnywhere, exposing an API.
  \item Frontend separated from the backend; Streamlit is the initial option.
  \item The frontend does not import the calculation modules.
  \item The backend does not query external data sources during user requests; it only reads the snapshot.
  \item The application domain uses object-oriented programming; other layers may use modular programming.
  \item Source code versioned in a Git repository.
  \item Errors handled with clear messages (invalid ticker, API failure, infeasible optimization).
  \item Source and extraction date of the data recorded.
  \item Raw data from B3, Yahoo and Bloomberg not published in the repository.
  \item Visible notice that the tool is academic and does not provide investment advice.
\end{itemize}

\section{Out of initial scope}
Items identified during elicitation that are left for future versions:
\begin{itemize}
  \item Automatic retry on API failures.
  \item Historical simulation of a user portfolio in the application.
  \item Alternative covariance estimators as defaults (Ledoit-Wolf may be used as a robustness check).
  \item Posterior covariance $\Sigma_{BL}$ in the optimization.
  \item User-defined view confidence, Idzorek's method and basket-versus-basket views.
  \item Free-float adjustment and the index theoretical portfolio as an alternative prior.
  \item Sensitivity analysis of $\tau$ and $\delta$, liquidity and outlier diagnostics.
  \item Bloomberg data, pending discussion with the course instructor.
\end{itemize}
